\documentclass[handout]{beamer}
\usetheme{Madrid}
\usecolortheme{default}

\usepackage[utf8]{inputenc}
\usepackage[ngerman]{babel}
\usepackage{listings}
\usepackage{xcolor}
\usepackage{tikz}
\usepackage{graphicx}

% Code listing style
\lstset{
    basicstyle=\ttfamily\footnotesize,
    keywordstyle=\color{blue}\bfseries,
    commentstyle=\color{green!60!black},
    stringstyle=\color{red},
    breaklines=true
}

\title{Einführung in die Kryptographie}
\subtitle{Von klassischen Chiffren zur modernen Verschlüsselung}
\author{Informatikkurs}
\date{\today}

\begin{document}

\frame{\titlepage}

\begin{frame}{Inhaltsverzeichnis}
\tableofcontents
\end{frame}

% =====================================================
\section{Einführung in die Kryptographie}
% =====================================================

\begin{frame}{Was ist Kryptographie?}
\begin{block}{Definition}
Die Wissenschaft und Praxis der Sicherung von Kommunikation durch Codes und Chiffren
\end{block}

\textbf{Etymologie:} Griechisch \emph{kryptos} (verborgen) + \emph{graphein} (schreiben) \pause

\vspace{0.5cm}
\textbf{Wir nutzen Kryptographie jeden Tag:}
\begin{itemize}
    \item<2-> Messaging-Apps (WhatsApp, Signal)
    \item<3-> Online-Shopping
    \item<4-> Banking
    \item<5-> WLAN-Verbindungen
\end{itemize}

% IMAGE SUGGESTION: Icon/diagram showing various everyday uses of cryptography
\end{frame}

\begin{frame}{Warum brauchen wir Kryptographie?}
\textbf{Ohne Kryptographie sind wir anfällig für:}

\begin{description}
    \item[Abhören] Unbefugte lesen deine Nachrichten
    \item[Manipulation] Angreifer verändern Nachrichten unbemerkt
    \item[Identitätsdiebstahl] Jemand gibt sich als du aus
    \item[Datendiebstahl] Persönliche Informationen werden gestohlen\pause
\end{description}


\vspace{0.5cm}
\textbf{Kryptographie bietet:}
\begin{itemize}
    \item \textbf{Vertraulichkeit:} Nur Befugte können lesen
    \item \textbf{Integrität:} Daten können nicht unbemerkt verändert werden
    \item \textbf{Authentifizierung:} Absenderidentität überprüfen
    \item \textbf{Nichtabstreitbarkeit:} Absender kann Senden nicht leugnen
\end{itemize}

% IMAGE SUGGESTION: Diagram showing attacker scenarios and cryptographic protections
\end{frame}

\begin{frame}{Grenzen der Kryptographie}


% IMAGE SUGGESTION: XKCD comic about voting software/security

\textbf{Selbst starke Verschlüsselung schützt nicht gegen:}
\begin{itemize}
    \item Social-Engineering-Angriffe
    \item Malware auf deinem Gerät
    \item Schlechte Passwortwahl
    \item Insider-Bedrohungen
    \item Physischen Zugang zu Geräten
    \item Implementierungsfehler
\end{itemize}\pause

\begin{alertblock}{Merke}
``Kryptographie ist wie ein wirklich gutes Schloss an deiner Tür. Es hilft nicht, wenn die Wände aus Papier sind.''
\end{alertblock}
\end{frame}

\begin{frame}{Grundlegende Begriffe}
\begin{description}
    \item[Klartext]<1-> Ursprüngliche, lesbare Nachricht
    \item[Geheimtext]<2-> Verschlüsselte, unlesbare Nachricht
    \item[Chiffre]<3-> Algorithmus zur Ver-/Entschlüsselung
    \item[Schlüssel]<4-> Geheime Information für die Chiffre
    \item[Verschlüsselung]<5-> Umwandlung von Klartext in Geheimtext
    \item[Entschlüsselung]<5-> Umwandlung von Geheimtext in Klartext
    \item[Schlüsselraum]<6-> Menge aller möglichen Schlüssel
\end{description}

% IMAGE SUGGESTION: Simple diagram showing plaintext → encryption → ciphertext → decryption → plaintext
\end{frame}

% \begin{frame}{Lernt Alice, Bob und Eve kennen}
% \textbf{Standardnamen in der Kryptographie:}

% \begin{description}
%     \item[Alice] Erste Partei, die eine Nachricht sendet
%     \item[Bob] Zweite Partei, die eine Nachricht empfängt
%     \item[Eve] Lauscherin, die versucht abzufangen
%     \item[Mallory] Böswilliger Angreifer, der Nachrichten verändert
%     \item[Trent] Vertrauenswürdige dritte Partei
% \end{description}

% % IMAGE SUGGESTION: Cartoon illustration of Alice, Bob, Eve showing communication flow
% \end{frame}

% =====================================================
\section{Klassische Kryptographie}
% =====================================================

\begin{frame}{Historischer Kontext}
\textbf{Kryptographie durch die Geschichte:}

\begin{itemize}
    \item \textbf{Altes Ägypten (1900 v. Chr.):} Hieroglyphen-Substitutionen
    \item \textbf{Spartaner (5. Jh. v. Chr.):} Skytale-Transpositionschiffre
    \item \textbf{Julius Cäsar (100-44 v. Chr.):} Cäsar-Chiffre für das Militär
    \item \textbf{Mittelalter:} Al-Kindis Häufigkeitsanalyse (850 n. Chr.)
    \item \textbf{2. Weltkrieg:} Enigma-Maschine und Bletchley Park
\end{itemize}

% IMAGE SUGGESTION: Timeline with historical cryptography examples/images
\end{frame}

\begin{frame}{Steganographie vs. Kryptographie}
\begin{block}{Steganographie}
Verbergen der \emph{Existenz} einer Nachricht
\end{block}

\textbf{Beispiele:}
\begin{itemize}
    \item Unsichtbare Tinte
    \item Mikropunkte
    \item Akrostichon
    \item Versteckte Daten in Bildern/Audio
\end{itemize}
\pause
\vspace{0.3cm}
\begin{exampleblock}{Probiere es aus}
\texttt{Geheim erscheint hier ein interessanter Markt,\\
nicht immer sichtbar.}
\end{exampleblock}\pause

Versteckte Nachricht: \textbf{GEHEIM} (erste Buchstaben)

% IMAGE SUGGESTION: Example of steganography - image with hidden data
\end{frame}

\begin{frame}[fragile]{Cäsar-Chiffre}
\textbf{Eine der ältesten Chiffren:}

Verschiebe jeden Buchstaben um eine feste Anzahl von Positionen

\vspace{0.3cm}
\textbf{Beispiel mit Verschiebung k=3:}
\begin{verbatim}
Klartext:   DER SCHNELLE FUCHS
Geheimtext: GHU VFKQHOOH IXFKV
\end{verbatim}

\vspace{0.3cm}
\textbf{Mathematische Notation:}
\begin{align*}
E(x) &= (x + k) \mod 26 \\
D(x) &= (x - k) \mod 26
\end{align*}

% IMAGE SUGGESTION: Caesar cipher wheel diagram
\end{frame}

\begin{frame}{Die Cäsar-Chiffre knacken}
\textbf{Warum die Cäsar-Chiffre schwach ist:}

\begin{itemize}
    \item \textbf{Kleiner Schlüsselraum:} Nur 25 mögliche Schlüssel
    \item \textbf{Brute-Force trivial:} Probiere alle Verschiebungen
    \item \textbf{Häufigkeitsanalyse funktioniert:} Häufigster Buchstabe wahrscheinlich = E
\end{itemize}

\vspace{0.5cm}
\begin{exampleblock}{Probiere es selbst}
Nutze \texttt{cryptii.com} um mit Cäsar-Chiffren zu experimentieren
\end{exampleblock}

% IMAGE SUGGESTION: Frequency analysis chart showing letter distribution
\end{frame}

\begin{frame}[fragile]{Monoalphabetische Substitution}
\textbf{Jede Permutation des Alphabets:}

\begin{verbatim}
Klar:     ABCDEFGHIJKLMNOPQRSTUVWXYZ
Geheim:   QWERTYUIOPASDFGHJKLZXCVBNM
\end{verbatim}

\textbf{Schlüsselraum:} $26! \approx 4 \times 10^{26}$ (riesig!)

\vspace{0.5cm}
\textbf{Aber immer noch anfällig für Häufigkeitsanalyse:}
\begin{itemize}
    \item Häufigste Buchstaben (DE): E, N, I, S, R, A, T
    \item Häufige Bigramme: EN, ER, CH, DE, EI
    \item Häufige Trigramme: EIN, ICH, DER
\end{itemize}

% IMAGE SUGGESTION: Frequency analysis bar chart for German text
\end{frame}

\begin{frame}[fragile]{Transpositionschiffren}
\textbf{Buchstaben umordnen statt ersetzen}

\vspace{0.3cm}
\textbf{Gartenzaun-Chiffre (3 Schienen):}

\begin{verbatim}
Nachricht: WIR SIND ENTDECKT FLIEHT SOFORT

W . . . I . . . N . . . C . . . L . . . T . . . O . .
. I . S . N . E . T . E . K . F . I . H . S . F . R .
. . R . . . D . . . E . . . T . . . E . . . O . . . T

Geheimtext: WINCL TOISE NTEKF IHSFR RDETT EOT
\end{verbatim}

% IMAGE SUGGESTION: Visual diagram of rail fence zigzag pattern
\end{frame}

\begin{frame}[fragile]{Vigenère-Chiffre}
\textbf{Polyalphabetische Substitution:} Verschiedene Verschiebungen an verschiedenen Positionen

\vspace{0.3cm}
\textbf{Beispiel mit Schlüsselwort LEMON:}
\begin{verbatim}
Klartext:   A N G R I F F A U F
Schlüssel:  L E M O N L E M O N
Geheimtext: L R S F V Q J M I S
\end{verbatim}

\vspace{0.3cm}
\begin{itemize}
    \item Genannt ``le chiffre indéchiffrable'' (die unknackbare Chiffre)
    \item Gleicher Klartextbuchstabe → verschiedene Geheimtextbuchstaben
    \item Überwindet einfache Häufigkeitsanalyse
\end{itemize}

% IMAGE SUGGESTION: Vigenère square (tabula recta)
\end{frame}

\begin{frame}{Vigenère knacken}
\textbf{Obwohl ``unknackbar'' genannt, kann sie gebrochen werden:}

\vspace{0.3cm}
\textbf{Methode:}
\begin{enumerate}
    \item Schlüssellänge finden (Kasiski-Untersuchung)
    \begin{itemize}
        \item Nach wiederholten Sequenzen suchen
        \item Abstände zwischen Wiederholungen teilen Faktoren mit Schlüssellänge
    \end{itemize}
    \item Geheimtext in Gruppen aufteilen (jeden $n$-ten Buchstaben)
    \item Jede Gruppe mit gleicher Cäsar-Verschiebung verschlüsselt
    \item Häufigkeitsanalyse auf jede Gruppe anwenden
\end{enumerate}
\end{frame}

\begin{frame}[fragile]{One-Time-Pad: Perfekte Geheimhaltung}
\textbf{Die einzige beweisbar unknackbare Chiffre!}

\vspace{0.3cm}
\textbf{Anforderungen:}
\begin{itemize}
    \item Schlüssel so lang wie die Nachricht
    \item Wirklich zufälliger Schlüssel
    \item Schlüssel nie wiederverwenden (auch nicht teilweise)
    \item Schlüssel nach Gebrauch vernichten
\end{itemize}
\pause
\vspace{0.3cm}
\begin{exampleblock}{Beispiel}
\begin{verbatim}
Klartext:   H A L L O
Schlüssel:  X M C K L (zufällig, gleiche Länge)
Geheimtext: E Q N V Z
\end{verbatim}
\end{exampleblock}

\textbf{Warum unknackbar:} Jeder Klartext gleich wahrscheinlich ohne Schlüssel!
\end{frame}

\begin{frame}{Warum wir das One-Time-Pad nicht nutzen}
\textbf{Praktische Herausforderungen:}

\begin{enumerate}
    \item \textbf{Schlüssellänge:} Muss der Nachrichtenlänge entsprechen (1 GB Video = 1 GB Schlüssel)
    \item \textbf{Schlüsselverteilung:} Wie teilt man riesige Zufallsschlüssel sicher?
    \item \textbf{Echte Zufälligkeit:} Schwer zu erzeugen
    \item \textbf{Schlüsselmanagement:} Nie wiederverwenden, muss vernichtet werden
    \item \textbf{Synchronisation:} Beide Parteien müssen gleichen Schlüsselteil nutzen
\end{enumerate}
\pause
\vspace{0.3cm}
\begin{alertblock}{Historisches Versagen}
Sowjetunion verwendete One-Time-Pads wieder → NSAs VENONA-Projekt entschlüsselte tausende Nachrichten
\end{alertblock}
\end{frame}

% =====================================================
\section{Moderne Kryptographie}
% =====================================================

\begin{frame}{Die kryptographische Revolution}
\begin{block}{Kerckhoffs' Prinzip (1883)}
``Ein Kryptosystem sollte sicher sein, selbst wenn alles über das System außer dem Schlüssel öffentlich bekannt ist.''
\end{block}

\vspace{0.3cm}
\textbf{Vorteile des modernen Ansatzes:}
\begin{itemize}
    \item Algorithmen können veröffentlicht und von Experten getestet werden
    \item Sicherheit beruht nur auf geheimen Schlüsseln
    \item Standardalgorithmen weltweit verwendet
    \item Bei Schlüsselkompromittierung nur Schlüssel wechseln (nicht ganzes System)
\end{itemize}
\end{frame}

\begin{frame}{Klassische vs. Moderne Kryptographie}
\begin{center}
\begin{tabular}{|l|l|l|}
\hline
\textbf{Aspekt} & \textbf{Klassisch} & \textbf{Modern} \\
\hline
Sicherheitsbasis & Geheimer Algorithmus & Geheimer Schlüssel \\
Schlüsselgröße & Klein & Groß (128-4096 Bits) \\
Grundlage & Einfache Mathematik & Komplexe Mathematik \\
Knackbarkeit & Musteranalyse & Berechnungshärte \\
Geschwindigkeit & Schnell (von Hand) & Benötigt Computer \\
Standardisierung & Schwierig & Weit verbreitet \\
\hline
\end{tabular}
\end{center}

% IMAGE SUGGESTION: Visual comparison diagram
\end{frame}

\begin{frame}{Berechnungssicherheit}
\textbf{Moderne Kryptographie beruht auf Berechnungshärte}

\vspace{0.3cm}
\textbf{Beispiel: 128-Bit-Schlüssel}
\begin{itemize}
    \item $2^{128} \approx 3,4 \times 10^{38}$ mögliche Werte
    \item 1 Billion Schlüssel pro Sekunde testen
    \item Würde $\approx 10^{21}$ Jahre dauern
    \item \emph{Weit länger als das Alter des Universums!}
\end{itemize}

\vspace{0.3cm}
\textbf{Kernaussage:}
Knacken ist theoretisch möglich, aber praktisch undurchführbar
\end{frame}

\begin{frame}{Drei Säulen der modernen Kryptographie}
\begin{enumerate}
    \item<1-> \textbf{Symmetrische Verschlüsselung}
    \begin{itemize}
        \item Gleicher Schlüssel für Ver-/Entschlüsselung
        \item Sehr schnell
        \item Beispiele: AES, ChaCha20
    \end{itemize}

    \item<2-> \textbf{Asymmetrische Verschlüsselung}
    \begin{itemize}
        \item Verschiedene Schlüssel (öffentlich/privat)
        \item Löst Schlüsselverteilung
        \item Beispiele: RSA, ECC
    \end{itemize}

    \item<3-> \textbf{Kryptographische Hashfunktionen}
    \begin{itemize}
        \item Einwegfunktionen
        \item Feste Ausgabegröße
        \item Beispiele: SHA-256, SHA-3
    \end{itemize}
\end{enumerate}

% IMAGE SUGGESTION: Diagram showing relationship between the three pillars
\end{frame}

% =====================================================
\section{Symmetrische Verschlüsselung}
% =====================================================

\begin{frame}{Überblick Symmetrische Verschlüsselung}
\textbf{Gleicher Schlüssel für beide Operationen}\pause

\vspace{0.3cm}
\textbf{Vorteile:}
\begin{itemize}
    \item Sehr schnell
    \item Starke Sicherheit mit kurzen Schlüsseln (256 Bits)
    \item Geringe Rechenanforderungen
\end{itemize}
\pause
\vspace{0.3cm}
\textbf{Nachteile:}
\begin{itemize}
    \item Schlüsselverteilungsproblem
    \item Schlüsselmanagement ($n$ Personen = $\frac{n(n-1)}{2}$ Schlüssel)
    \item Schlüsselkompromittierung betrifft alle Nachrichten
\end{itemize}

% IMAGE SUGGESTION: Diagram showing Alice and Bob with shared key
\end{frame}

\begin{frame}{Blockchiffren vs. Stromchiffren}
\begin{columns}
\column{0.5\textwidth}
\textbf{Blockchiffren}
\begin{itemize}
    \item Feste Blockgröße (128 Bits)
    \item Beispiele: AES, DES
    \item Brauchen Padding
    \item Benötigen Betriebsmodus
\end{itemize}
\pause

\column{0.5\textwidth}
\textbf{Stromchiffren}
\begin{itemize}
    \item Bit für Bit oder Byte für Byte
    \item Beispiele: ChaCha20, RC4
    \item Kein Padding nötig
    \item Gut für Streaming-Daten
\end{itemize}
\end{columns}

\vspace{0.5cm}
% IMAGE SUGGESTION: Visual comparison of block vs stream cipher operation
\end{frame}

\begin{frame}{Betriebsmodi für Blockchiffren}
\textbf{Problem:} Identische Klartextblöcke → identische Geheimtextblöcke

\vspace{0.3cm}
\textbf{Gängige Modi:}
\begin{description}
    \item[ECB] Electronic Codebook - \alert{NICHT SICHER!}
    \item[CBC] Cipher Block Chaining - jeder Block mit vorherigem XOR-verknüpft
    \item[CTR] Counter Mode - macht Blockchiffre zur Stromchiffre
    \item[GCM] Galois/Counter Mode - Verschlüsselung + Authentifizierung
\end{description}

% IMAGE SUGGESTION: ECB Penguin example showing why ECB is insecure
\end{frame}

\begin{frame}{AES: Advanced Encryption Standard}
\textbf{Meistgenutzte symmetrische Verschlüsselung (eingeführt 2001)}

\vspace{0.3cm}
\textbf{Spezifikationen:}
\begin{itemize}
    \item Blockgröße: 128 Bits
    \item Schlüsselgröße: 128, 192 oder 256 Bits
    \item Runden: 10, 12 oder 14 (je nach Schlüsselgröße)
    \item Basiert auf Substitutions-Permutations-Netzwerk
\end{itemize}

\vspace{0.3cm}
\textbf{Überall verwendet:}
\begin{itemize}
    \item WLAN (WPA2/WPA3)
    \item VPNs und TLS/SSL
    \item Dateiverschlüsselung
    \item Messaging-Apps
    \item Regierung TOP SECRET (mit 256-Bit-Schlüsseln)
\end{itemize}

% IMAGE SUGGESTION: AES encryption round diagram
\end{frame}

\begin{frame}{AES-Sicherheit}
\textbf{AES-256 hat $2^{256}$ mögliche Schlüssel}

\vspace{0.3cm}
Bei 1 Billion Billionen Schlüsseln pro Sekunde testen:
\[
\frac{2^{256}}{10^{24}} \approx 3,7 \times 10^{51} \text{ Sekunden} \approx 10^{44} \text{ Jahre}
\]

\vspace{0.3cm}
\begin{alertblock}{Kernaussage}
Verdopplung der Schlüssellänge verdoppelt nicht die Sicherheit - sie \emph{quadriert} sie!

$2^{256} = (2^{128})^2$ bedeutet 256-Bit ist $(2^{128})$ mal sicherer als 128-Bit
\end{alertblock}
\end{frame}

% =====================================================
\section{Asymmetrische Verschlüsselung}
% =====================================================

\begin{frame}{Das Schlüsselverteilungsproblem}
\textbf{Herausforderung bei symmetrischer Verschlüsselung:}

\begin{itemize}
    \item Alice und Bob brauchen gemeinsamen geheimen Schlüssel
    \item Haben sich nie persönlich getroffen
    \item Jeder über Internet gesendete Schlüssel könnte abgefangen werden
    \item Kann den Schlüssel nicht verschlüsseln (braucht anderen Schlüssel - unendlicher Regress!)
\end{itemize}

\vspace{0.5cm}
\textbf{Dieses Problem plagte die Kryptographie jahrhundertelang...}

% IMAGE SUGGESTION: Diagram showing Alice, Bob, and Eve with intercepted key
\end{frame}

\begin{frame}{Die revolutionäre Idee}
\textbf{Public-Key-Kryptographie (Diffie-Hellman, 1976)}

\vspace{0.3cm}
\textbf{Kernkonzept:}
\begin{itemize}
    \item Zwei mathematisch verbundene Schlüssel: \textbf{öffentlich} und \textbf{privat}
    \item Öffentlicher Schlüssel frei geteilt
    \item Privater Schlüssel geheim gehalten
    \item Mit öffentlichem verschlüsselt → mit privatem entschlüsselt
    \item Berechnungstechnisch unmöglich, privaten aus öffentlichem abzuleiten
\end{itemize}

\vspace{0.3cm}
\begin{block}{Briefkasten-Analogie}
Öffentlicher Schlüssel = Briefkasten (jeder kann Briefe einwerfen)\\
Privater Schlüssel = Briefkastenschlüssel (nur du kannst öffnen)
\end{block}

% IMAGE SUGGESTION: Mailbox analogy illustration
\end{frame}

\begin{frame}{Wie es die Schlüsselverteilung löst}
\begin{enumerate}
    \item Bob erzeugt öffentlich-privates Schlüsselpaar
    \item Bob veröffentlicht öffentlichen Schlüssel
    \item Alice schlägt Bobs öffentlichen Schlüssel nach
    \item Alice verschlüsselt Nachricht mit Bobs öffentlichem Schlüssel
    \item Alice sendet Geheimtext
    \item Bob entschlüsselt mit seinem privaten Schlüssel
\end{enumerate}

\vspace{0.3cm}
\textbf{Keine geheimen Informationen vorher ausgetauscht!}

% IMAGE SUGGESTION: Step-by-step diagram of public key encryption process
\end{frame}

\begin{frame}{Symmetrisch vs. Asymmetrisch}
\begin{center}
\begin{tabular}{|l|l|l|}
\hline
\textbf{Eigenschaft} & \textbf{Symmetrisch} & \textbf{Asymmetrisch} \\
\hline
Schlüssel & Gleicher Schlüssel & Öffentlich + Privat \\
Schlüsselverteilung & Schwierig & Einfach \\
Geschwindigkeit & Sehr schnell & Viel langsamer \\
Schlüsselgröße & 128-256 Bits & 2048-4096 Bits \\
Verwendung & Massenverschlüsselung & Schlüsselaustausch \\
Beispiele & AES & RSA, ECC \\
\hline
\end{tabular}
\end{center}

\vspace{0.3cm}
\textbf{In der Praxis:} Kombiniere beide für beste Ergebnisse!
\end{frame}

\begin{frame}{RSA-Verschlüsselung}
\textbf{RSA (Rivest-Shamir-Adleman, 1977)}

Basiert auf der Schwierigkeit, große Zahlen zu faktorisieren

\vspace{0.3cm}
\textbf{Schlüsselerzeugung:}
\begin{enumerate}
    \item Wähle zwei große Primzahlen $p$ und $q$
    \item Berechne $n = p \times q$ (öffentlich)
    \item Berechne $\phi(n) = (p-1)(q-1)$
    \item Wähle öffentlichen Exponenten $e$ (üblicherweise 65537)
    \item Berechne privaten Exponenten $d$
    \item Öffentlicher Schlüssel: $(n, e)$, Privater Schlüssel: $(n, d)$
\end{enumerate}

\vspace{0.3cm}
\textbf{Verschlüsselung:} $\text{Geheimtext} = \text{Nachricht}^e \mod n$

\textbf{Entschlüsselung:} $\text{Nachricht} = \text{Geheimtext}^d \mod n$

% IMAGE SUGGESTION: RSA key generation flowchart
\end{frame}

\begin{frame}{RSA in der Praxis}
\textbf{Schlüsselgrößen:}
\begin{itemize}
    \item 1024 Bits: \alert{Veraltet} (von Nationalstaaten faktorisiert)
    \item 2048 Bits: Aktuelles Minimum
    \item 3072 Bits: Hohe Sicherheit
    \item 4096 Bits: Sehr hohe Sicherheit (langsamer)
\end{itemize}

\vspace{0.3cm}
\textbf{Häufige Verwendungen:}
\begin{itemize}
    \item TLS/SSL-Zertifikate (HTTPS)
    \item E-Mail-Verschlüsselung (PGP/GPG)
    \item SSH-Schlüssel
    \item Digitale Signaturen
\end{itemize}

\vspace{0.3cm}
\textbf{Hinweis:} Zu langsam für große Daten - hybride Verschlüsselung nutzen!
\end{frame}

\begin{frame}{Elliptische-Kurven-Kryptographie (ECC)}
\textbf{Gleiche Sicherheit wie RSA mit viel kleineren Schlüsseln}

\vspace{0.3cm}
\begin{center}
\begin{tabular}{|c|c|}
\hline
\textbf{RSA-Schlüsselgröße} & \textbf{ECC-Schlüsselgröße} \\
\hline
1024 Bits & 160 Bits \\
2048 Bits & 224 Bits \\
3072 Bits & 256 Bits \\
15360 Bits & 512 Bits \\
\hline
\end{tabular}
\end{center}

\vspace{0.3cm}
\textbf{Vorteile:}
\begin{itemize}
    \item Kleinere Schlüssel, Zertifikate, Signaturen
    \item Schnellere Operationen
    \item Weniger Bandbreite und Speicher
\end{itemize}

% IMAGE SUGGESTION: Elliptic curve graph
\end{frame}

\begin{frame}{ECC-Anwendungen}
\textbf{Wo ECC verwendet wird:}
\begin{itemize}
    \item Bitcoin und Kryptowährungs-Wallets
    \item Moderne TLS/SSL-Zertifikate
    \item Signal, WhatsApp (Curve25519)
    \item Apple iMessage
    \item SSH-Schlüssel
    \item IoT-Geräte (begrenzte Rechenleistung)
\end{itemize}

\vspace{0.3cm}
\textbf{Beliebte Kurven:}
\begin{itemize}
    \item P-256, P-384, P-521 (NIST-Standard)
    \item Curve25519 (modern, sicher)
    \item Ed25519 (optimiert für Signaturen)
\end{itemize}
\end{frame}

% =====================================================
\section{Kryptographische Hashfunktionen}
% =====================================================

\begin{frame}[fragile]{Was ist eine Hashfunktion?}
\textbf{Nimmt beliebige Eingabe → erzeugt Ausgabe fester Größe}

\vspace{0.3cm}
\begin{exampleblock}{Beispiel: SHA-256}
\texttt{Eingabe: "Hallo, Welt!"}\\
\texttt{Hash: dffd6021bb2bd5b0af676290809ec3a5...}

\vspace{0.2cm}
\texttt{Eingabe: "Hallo, Welt?" (beachte das ?)}\\
\texttt{Hash: a0e0d1d23f58c99b8a73fdecd1ba0d8f...}
\end{exampleblock}

\vspace{0.3cm}
\textbf{Beachte:}
\begin{itemize}
    \item Gleiche Ausgabelänge (256 Bits)
    \item Winzige Änderung → völlig anderer Hash
    \item Sieht zufällig aus
\end{itemize}

% IMAGE SUGGESTION: Visual diagram of hash function with avalanche effect
\end{frame}

\begin{frame}{Eigenschaften kryptographischer Hashes}
\begin{enumerate}
    \item \textbf{Deterministisch:} Gleiche Eingabe → gleicher Hash
    \item \textbf{Schnell zu berechnen}
    \item \textbf{Urbildresistenz:} Schwer, Eingabe aus Hash zu finden
    \item \textbf{Zweite-Urbild-Resistenz:} Schwer, andere Eingabe mit gleichem Hash zu finden
    \item \textbf{Kollisionsresistenz:} Schwer, zwei beliebige Eingaben mit gleichem Hash zu finden
    \item \textbf{Lawineneffekt:} Kleine Änderung → drastisch anderer Hash
    \item \textbf{Feste Ausgabegröße}
\end{enumerate}
\end{frame}

\begin{frame}{Gängige Hashfunktionen}
\begin{center}
\begin{tabular}{|l|c|c|}
\hline
\textbf{Algorithmus} & \textbf{Ausgabegröße} & \textbf{Status} \\
\hline
MD5 & 128 Bits & \alert{Gebrochen} \\
SHA-1 & 160 Bits & \alert{Veraltet} \\
SHA-256 & 256 Bits & \textcolor{green}{Sicher} \\
SHA-512 & 512 Bits & \textcolor{green}{Sicher} \\
SHA-3 & Variabel & \textcolor{green}{Sicher} \\
BLAKE2/3 & Variabel & \textcolor{green}{Sicher} \\
\hline
\end{tabular}
\end{center}

\vspace{0.3cm}
\begin{alertblock}{Wichtig}
Niemals MD5 oder SHA-1 für Sicherheitszwecke verwenden!
\end{alertblock}
\end{frame}

\begin{frame}{Anwendungen von Hashfunktionen}
\textbf{1. Datenintegrität / Prüfsummen}
\begin{itemize}
    \item Datei-Downloads verifizieren
    \item Korruption oder Manipulation erkennen
\end{itemize}

\vspace{0.3cm}
\textbf{2. Passwortspeicherung}
\begin{itemize}
    \item Niemals Klartext-Passwörter speichern!
    \item Hash des Passworts speichern
    \item Salt hinzufügen gegen Rainbow-Tables
    \item Spezialisierte Funktionen nutzen: Argon2, bcrypt
\end{itemize}

\vspace{0.3cm}
\textbf{3. Digitale Signaturen}
\begin{itemize}
    \item Dokument hashen, dann Hash signieren
\end{itemize}

% IMAGE SUGGESTION: Diagram showing password hashing with salt
\end{frame}

\begin{frame}[fragile]{Passwort-Hashing mit Salt}
\textbf{Salt = zufälliger String, der vor dem Hashen hinzugefügt wird}

\vspace{0.3cm}
\begin{exampleblock}{Beispiel}
Benutzer 1: "passwort123" + Salt "a8f9e2" → Hash 1\\
Benutzer 2: "passwort123" + Salt "x3k8m1" → Hash 2
\end{exampleblock}

\vspace{0.3cm}
\textbf{In Datenbank gespeichert:}
\begin{center}
\begin{tabular}{|l|l|l|}
\hline
\textbf{Benutzer} & \textbf{Salt} & \textbf{Hash} \\
\hline
alice & a8f9e2 & 3f7a8c9d... \\
bob & x3k8m1 & 9e2b5f1c... \\
\hline
\end{tabular}
\end{center}

\textbf{Verhindert:}
\begin{itemize}
    \item Rainbow-Table-Angriffe
    \item Identifizierung von Benutzern mit gleichem Passwort
    \item Parallele Angriffe
\end{itemize}
\end{frame}

\begin{frame}{Blockchain und Kryptowährungen}
\textbf{Hashfunktionen sind fundamental für Blockchain:}

\vspace{0.3cm}
\begin{itemize}
    \item \textbf{Block-Hashing:} Jeder Block enthält Hash des vorherigen (erzeugt Kette)
    \item \textbf{Proof of Work:} Finde Nonce, bei der Block-Hash mit Nullen beginnt
    \item \textbf{Adressen:} Bitcoin-Adressen sind Hashes öffentlicher Schlüssel
    \item \textbf{Merkle-Bäume:} Effiziente Transaktionsverifizierung
\end{itemize}

% IMAGE SUGGESTION: Blockchain diagram showing hash links between blocks
\end{frame}

% =====================================================
\section{Digitale Signaturen}
% =====================================================

\begin{frame}{Digitale Signaturen}
\textbf{Elektronisches Äquivalent der handschriftlichen Unterschrift}

\vspace{0.3cm}
\textbf{Bietet:}
\begin{itemize}
    \item \textbf{Authentifizierung:} Beweist, wer die Nachricht gesendet hat
    \item \textbf{Integrität:} Beweist, dass Nachricht nicht verändert wurde
    \item \textbf{Nichtabstreitbarkeit:} Absender kann Senden nicht leugnen
\end{itemize}

\vspace{0.3cm}
\textbf{Wie es funktioniert:}
\begin{enumerate}
    \item Nachricht hashen
    \item Hash mit \textbf{privatem Schlüssel} verschlüsseln (Signatur)
    \item Nachricht + Signatur senden
    \item Empfänger entschlüsselt Signatur mit \textbf{öffentlichem Schlüssel}
    \item Mit Hash der empfangenen Nachricht vergleichen
\end{enumerate}

% IMAGE SUGGESTION: Digital signature process diagram
\end{frame}

\begin{frame}{Verschlüsselung vs. Signieren}
\begin{center}
\begin{tabular}{|l|l|l|}
\hline
\textbf{Eigenschaft} & \textbf{Verschlüsselung} & \textbf{Signieren} \\
\hline
Sperren mit & Öffentlichem Schlüssel & Privatem Schlüssel \\
Entsperren mit & Privatem Schlüssel & Öffentlichem Schlüssel \\
Zweck & Vertraulichkeit & Authentifizierung \\
Wer kann es? & Jeder & Nur Schlüsselbesitzer \\
Wer kann verifizieren? & Nur Schlüsselbesitzer & Jeder \\
\hline
\end{tabular}
\end{center}

\vspace{0.5cm}
\textbf{Merke:}
\begin{itemize}
    \item \textbf{Verschlüsselung:} Öffentlich sperren, privat entsperren (Geheimhaltung)
    \item \textbf{Signieren:} Privat sperren, öffentlich entsperren (Authentizität)
\end{itemize}
\end{frame}

\begin{frame}{Public-Key-Infrastruktur (PKI)}
\textbf{Problem:} Woher weißt du, dass ein öffentlicher Schlüssel dem gehört, der er behauptet?

\vspace{0.3cm}
\textbf{Lösung: Digitale Zertifikate}

\vspace{0.3cm}
\textbf{Zertifikat enthält:}
\begin{itemize}
    \item Öffentlichen Schlüssel
    \item Identität des Besitzers
    \item Gültigkeitszeitraum
    \item Digitale Signatur der Zertifizierungsstelle (CA)
\end{itemize}

\vspace{0.3cm}
\textbf{Vertrauenskette:}
Website → Zwischen-CA → Root-CA (im Browser vorinstalliert)

% IMAGE SUGGESTION: PKI certificate chain diagram
\end{frame}

\begin{frame}{Wie HTTPS funktioniert}
\textbf{Wenn du \texttt{https://www.example.com} besuchst:}

\begin{enumerate}
    \item Server sendet Zertifikat (mit öffentlichem Schlüssel)
    \item Browser verifiziert:
    \begin{itemize}
        \item Zertifikat von vertrauenswürdiger CA signiert
        \item Nicht abgelaufen
        \item Stimmt mit Domainnamen überein
        \item Nicht widerrufen
    \end{itemize}
    \item Browser und Server führen Schlüsselaustausch durch
    \item Etablieren symmetrischen Verschlüsselungsschlüssel
    \item Alle Daten mit Sitzungsschlüssel verschlüsselt
    \item Grünes Schloss erscheint!
\end{enumerate}

\vspace{0.3cm}
\textbf{Bietet:} Vertraulichkeit + Authentifizierung + Integrität

% IMAGE SUGGESTION: HTTPS handshake diagram
\end{frame}

% =====================================================
\section{Moderne Anwendungen}
% =====================================================

\begin{frame}{Ende-zu-Ende-Verschlüsselung (E2EE)}
\textbf{Nur Sender und Empfänger können Nachrichten lesen}

Nicht einmal der Dienstanbieter!

\vspace{0.3cm}
\textbf{Beispiele:}
\begin{itemize}
    \item Signal
    \item WhatsApp
    \item iMessage
    \item ProtonMail
\end{itemize}

\vspace{0.3cm}
\textbf{Im Gegensatz zu:}
\begin{itemize}
    \item Normale E-Mail: Anbieter kann lesen
    \item SMS: Mobilfunkanbieter kann lesen
    \item Social-Media-DMs: Plattform kann lesen
\end{itemize}

% IMAGE SUGGESTION: E2EE vs non-E2EE comparison diagram
\end{frame}

\begin{frame}{Passwort-Manager}
\textbf{Speichern alle Passwörter sicher verschlüsselt}

\vspace{0.3cm}
\textbf{Wie sie funktionieren:}
\begin{itemize}
    \item Master-Passwort leitet Verschlüsselungsschlüssel ab
    \item Alle Passwörter mit AES-256 verschlüsselt
    \item Nur du kennst das Master-Passwort
\end{itemize}

\vspace{0.3cm}
\textbf{Warum einen nutzen:}
\begin{itemize}
    \item Einzigartiges starkes Passwort für jede Seite
    \item Schutz vor Phishing
    \item Schutz vor Keyloggern
    \item Keine Notwendigkeit, Passwörter zu merken
\end{itemize}

\vspace{0.3cm}
\textbf{Beispiele:} 1Password, Bitwarden, KeePass
\end{frame}

\begin{frame}{VPNs (Virtuelle Private Netzwerke)}
\textbf{Erstellen verschlüsselte Tunnel für Internetverkehr}

\vspace{0.3cm}
\textbf{Was VPNs tun:}
\begin{itemize}
    \item Verschlüsseln Verkehr zwischen dir und VPN-Server
    \item Verbergen deine IP-Adresse
    \item Umgehen geografische Beschränkungen
    \item Schützen in öffentlichen WLANs
\end{itemize}

\vspace{0.3cm}
\textbf{Was VPNs nicht tun:}
\begin{itemize}
    \item Dich komplett anonym machen
    \item Vor Malware schützen
    \item Daten nach dem VPN-Server verschlüsseln
\end{itemize}

% IMAGE SUGGESTION: VPN tunnel diagram
\end{frame}

% =====================================================
\section{Bedrohungen und Zukunft}
% =====================================================

\begin{frame}{Häufige Angriffe}
\begin{description}
    \item[Brute Force] Jeden möglichen Schlüssel ausprobieren
    \item[Wörterbuchangriff] Häufige Passwörter ausprobieren
    \item[Rainbow Tables] Vorberechnete Passwort-Hashes
    \item[Man-in-the-Middle] Kommunikation abfangen
    \item[Seitenkanalangriffe] Geheimnisse über Timing, Stromverbrauch extrahieren
    \item[Implementierungsfehler] Heartbleed, POODLE usw.
\end{description}

\vspace{0.3cm}
\textbf{Verteidigung:} Große Schlüssel, Salting, Authentifizierung, konstante Zeitalgorithmen, Updates
\end{frame}

\begin{frame}{Die Quantencomputer-Bedrohung}
\textbf{Quantencomputer könnten aktuelle Public-Key-Krypto brechen}

\vspace{0.3cm}
\textbf{Shors Algorithmus (1994):}
\begin{itemize}
    \item Kann große Zahlen effizient faktorisieren
    \item Bricht: RSA, Diffie-Hellman, ECC
\end{itemize}

\vspace{0.3cm}
\textbf{Symmetrische Verschlüsselung weniger betroffen:}
\begin{itemize}
    \item Grovers Algorithmus nur quadratische Beschleunigung
    \item Verdopplung der Schlüsselgröße stellt Sicherheit wieder her
\end{itemize}

\vspace{0.3cm}
\textbf{Zeitrahmen:} 10-30+ Jahre bis die Bedrohung real ist

\textbf{Lösung:} Post-Quanten-Kryptographie (PQC)

% IMAGE SUGGESTION: Timeline showing quantum computing development
\end{frame}

\begin{frame}{Post-Quanten-Kryptographie}
\textbf{NIST standardisierte quantenresistente Algorithmen:}

\vspace{0.3cm}
\begin{itemize}
    \item \textbf{CRYSTALS-Kyber} (Verschlüsselung)
    \item \textbf{CRYSTALS-Dilithium} (Signaturen)
    \item \textbf{FALCON} (Signaturen)
    \item \textbf{SPHINCS+} (Signaturen)
\end{itemize}

\vspace{0.3cm}
\textbf{Basieren auf Problemen, die Quantencomputer nicht lösen können:}
\begin{itemize}
    \item Gitterprobleme
    \item Hash-basierte Schemata
    \item Multivariate Polynome
\end{itemize}
\end{frame}

\begin{frame}{Social Engineering}
\begin{alertblock}{Merke}
``Das schwächste Glied ist oft der Mensch, nicht die Kryptographie.''
\end{alertblock}

\vspace{0.3cm}
\textbf{Häufige Taktiken:}
\begin{itemize}
    \item Phishing-E-Mails
    \item Vorwand (Identitätsdiebstahl)
    \item Ködern (infizierte USB-Sticks)
    \item Tailgating (physischer Zugang)
\end{itemize}

\vspace{0.3cm}
\textbf{Verteidigung:}
\begin{itemize}
    \item Sicherheitsbewusstseinsschulungen
    \item Zwei-Faktor-Authentifizierung
    \item Prinzip der minimalen Rechte
\end{itemize}

% IMAGE SUGGESTION: Phishing example or social engineering attack diagram
\end{frame}

% =====================================================
\section{Best Practices}
% =====================================================

\begin{frame}{Best Practices für Benutzer}
\begin{enumerate}
    \item Verwende starke, einzigartige Passwörter (12+ Zeichen)
    \item Verwende einen Passwort-Manager
    \item Aktiviere Zwei-Faktor-Authentifizierung (2FA)
    \item Halte Software aktuell
    \item Nutze HTTPS (achte auf das Schloss)
    \item Sei vorsichtig bei Phishing
    \item Nutze Ende-zu-Ende-verschlüsselte Messenger
    \item Verschlüssele sensible Daten (Festplattenverschlüsselung)
    \item Nutze VPN in öffentlichen WLANs
\end{enumerate}
\end{frame}

\begin{frame}{Best Practices für Entwickler}
\begin{enumerate}
    \item \textbf{Entwickle keine eigene Krypto!}
    \item Nutze etablierte Bibliotheken (NaCl, OpenSSL)
    \item Nutze moderne Algorithmen (AES-256, SHA-256, Argon2)
    \item Speichere niemals Klartext-Passwörter
    \item Nutze sichere Zufallszahlengeneratoren
    \item Implementiere Ratenbegrenzung
    \item Nutze TLS 1.3 mit starken Chiffren
    \item Halte Abhängigkeiten aktuell
    \item Führe Sicherheitsaudits durch
\end{enumerate}
\end{frame}

\begin{frame}{Häufige Fehler vermeiden}
\begin{itemize}
    \item Veraltete Algorithmen verwenden (MD5, SHA-1, DES)
    \item Passwörter im Klartext speichern
    \item Schlüssel, IVs oder Nonces wiederverwenden
    \item ECB-Modus verwenden
    \item Eigene Krypto implementieren
    \item Benutzereingaben vertrauen
    \item Timing-Angriffe ignorieren
    \item Schlüssel im Quellcode hartcodieren
    \item Schwache Zufallsgeneratoren verwenden
\end{itemize}
\end{frame}

% =====================================================
\section{Fazit}
% =====================================================

\begin{frame}{Wichtige Erkenntnisse}
\begin{enumerate}
    \item \textbf{Kryptographie ist essenziell} für unser digitales Leben
    \item \textbf{Sicherheit ist schwer} - kleine Fehler = große Konsequenzen
    \item \textbf{Nutze etablierte Tools} - implementiere keine eigenen
    \item \textbf{Verteidigung in der Tiefe} - Krypto ist notwendig, aber nicht ausreichend
    \item \textbf{Bleib informiert} - Bedrohungen entwickeln sich ständig
    \item \textbf{Der menschliche Faktor zählt} - stärkste Krypto kann Phishing nicht stoppen
\end{enumerate}
\end{frame}

\begin{frame}{Die Reise}
\textbf{Von antiken Chiffren zu modernen Systemen:}

\begin{itemize}
    \item Klassisch: Cäsar, Vigenère, One-Time-Pad
    \item Symmetrisch: AES, ChaCha20
    \item Asymmetrisch: RSA, ECC, Diffie-Hellman
    \item Hashfunktionen: SHA-256, Argon2
    \item Anwendungen: HTTPS, E2EE, Blockchain
    \item Zukunft: Post-Quanten-Kryptographie
\end{itemize}

\vspace{0.5cm}
\textbf{Kryptographie entwickelt sich weiter, um neuen Herausforderungen zu begegnen!}
\end{frame}

\begin{frame}{Weitere Ressourcen}
\textbf{Bücher:}
\begin{itemize}
    \item \emph{The Code Book} von Simon Singh
    \item \emph{Applied Cryptography} von Bruce Schneier
    \item \emph{Cryptography Engineering} von Ferguson et al.
\end{itemize}

\vspace{0.3cm}
\textbf{Online:}
\begin{itemize}
    \item cryptohack.org - Lernen durch Herausforderungen
    \item cryptopals.com - Praktische Übungen
    \item cryptii.com - Mit Chiffren experimentieren
    \item Coursera: Stanford Kryptographie-Kurs
\end{itemize}
\end{frame}

\begin{frame}
\begin{center}
\Huge Vielen Dank!

\vspace{1cm}
\large Fragen?

\vspace{1cm}
\normalsize
\textit{``Mit großer Macht kommt große Verantwortung.''}

\vspace{0.5cm}
Nutze Kryptographie, um Privatsphäre und Sicherheit zu schützen, nicht um Schaden anzurichten.
\end{center}
\end{frame}

\end{document}
